% !TeX program = lualatex
\documentclass[11pt,a4paper]{article}

% ========= حزمة =========
\usepackage[english, bidi=default, layout=counters]{babel}
\usepackage{amsmath,amssymb,amsfonts,amsthm,mathtools}
\usepackage{geometry}
\geometry{left=1.25cm,right=1.25cm,top=1.25cm,bottom=3.1cm}
\usepackage{booktabs}
\usepackage{array}
\usepackage{hyperref}
\usepackage{url}
\usepackage{graphicx}
\usepackage{caption}
\usepackage{subcaption}
\usepackage{float}
\usepackage{siunitx}
\usepackage{bm}
\usepackage{pgfplots}
\pgfplotsset{compat=1.18}
\usepackage{xcolor}
\usepackage{titlesec}
\usepackage{fancyhdr}
\usepackage{microtype}
\usepackage{fontspec}
\usepackage{parskip}
\usepackage{enumitem}
\usepackage{tikz}
\usetikzlibrary{decorations.pathreplacing,arrows.meta,positioning,shapes.geometric}

% ========= اللغات =========
\babelprovide[import, main, maparabic]{arabic}
\babelprovide[import, onchar=ids fonts]{english}

% ========= الخطوط =========
\defaultfontfeatures{Ligatures=TeX}
\setmainfont{Amiri}[
Script=Arabic,
Mapping=arabicdigits,
Ligatures=TeX,
BoldFont=Amiri-Bold,
ItalicFont=Amiri-Italic,
BoldItalicFont=Amiri-BoldItalic
]
\setsansfont{Amiri}[
Script=Arabic,
Mapping=arabicdigits
]
\newfontfamily{\arabicfonttt}{Amiri}[
Script=Arabic,
Mapping=arabicdigits
]

% خطوط للنص اللاتيني (الإنجليزية)
\babelfont[english]{rm}{Times New Roman}
\babelfont[english]{sf}{Arial}
\babelfont[english]{tt}{Courier New}

% ========= ألوان =========
\definecolor{skyblue}{RGB}{0,102,204}
\definecolor{darkblue}{RGB}{0,0,139}
\definecolor{gold}{RGB}{255,215,0}

% ========= أوامر YUCT =========
\newcommand{\Keff}{K_{\mathrm{eff}}}
\DeclareMathOperator{\Tr}{Tr}

% ========= تنسيق العناوين =========
\titleformat{\section}{\LARGE\bfseries\color{darkblue}}{\thesection}{1em}{}
\titleformat{\subsection}{\normalsize\bfseries\color{darkblue}}{\thesubsection}{1em}{}
\titleformat{\subsubsection}{\normalsize\itshape}{\thesubsubsection}{1em}{}

% ========= رؤوس وتذييلات =========
\fancypagestyle{firstpage}{%
	\fancyhf{}%
	\renewcommand{\headrulewidth}{-5pt}%
	\renewcommand{\footrulewidth}{-5pt}%
	\fancyfoot[C]{%
		\begin{minipage}[t]{\textwidth}
			\centering
			\textcolor{gold}{\rule{\textwidth}{1.6pt}}\\[5pt]
			{\small \textsuperscript{1}أبحاث ياكوشيف، \url{https://yuct.org/}} \hfill
			{\small بروتوكول YPSDC: \url{https://ypsdc.com/}}\\[-2pt]
			\textcolor{gold}{\rule{\textwidth}{1.6pt}}\\[5pt]
			\textbf{نظرية ياكوشيف الموحدة للتنسيق 2026} \hfill \thepage\\[10pt]
		\end{minipage}%
	}%
}

\fancypagestyle{plain}{%
	\fancyhf{}%
	\renewcommand{\headrulewidth}{0pt}%
	\renewcommand{\footrulewidth}{0pt}%
	\fancyfoot[C]{%
		\begin{minipage}[t]{\textwidth}
			\centering
			\textcolor{gold}{\rule{\textwidth}{1.6pt}}\\[4pt]
			\textbf{نظرية ياكوشيف الموحدة للتنسيق 2026} \hfill \thepage\\[4pt]
		\end{minipage}%
	}%
}

\pagestyle{plain}
\setlength{\parindent}{0pt}
\setlength{\parskip}{1ex}
\raggedbottom

\hypersetup{
	colorlinks=true,
	linkcolor=blue,
	citecolor=blue,
	urlcolor=blue,
}

% ========= رأس المقال (YUCT logo) =========
\newcommand{\makeheader}{%
	\noindent
	\begin{minipage}{0.17\textwidth}
		\raggedright
		{\sffamily\bfseries\color{skyblue}\fontsize{46}{46}\selectfont YUCT}%
	\end{minipage}%
	\hspace{30pt}
	\begin{minipage}{0.32\textwidth}
		\raggedright
		{\sffamily\fontsize{11}{13}\selectfont نظرية ياكوشيف\\ الموحدة\\ للتنسيق}%
	\end{minipage}%
	\hfill
	\begin{minipage}{0.38\textwidth}
		\raggedleft
		{\sffamily\bfseries مذكرة}\\
		{\sffamily مايو 2026}\\
		{\sffamily\footnotesize \url{https://yuct.org/}}%
	\end{minipage}
	\par\vspace{5pt}
	\noindent\textcolor{gold}{\rule{\textwidth}{1.6pt}}
	\par\vspace{0pt}
}

\begin{document}
	\renewcommand{\thepage}{\arabic{page}}
	\thispagestyle{firstpage}
	\makeheader
	
	\vspace{0cm}
	{\LARGE\bfseries YPSDC اللامركزي (d‑YPSDC): \\ التنسيق بدون تواصل مباشر عبر مؤشر بيئي مشترك \par}
	\vspace{0.2cm}
	
	{\large أليكسي ف. ياكوشيف\footnote{أبحاث ياكوشيف، \url{https://yuct.org/}}} \\
	\vspace{0.0cm}
	
	{\color{darkblue}\textbf{\LARGE المستخلص}} \par
	{\large
		\noindent
		يفترض بروتوكول YPSDC الكلاسيكي طوبولوجيا مركزية: وكيل واحد يرسل بنشاط مؤشراً قصيراً إلى آخر، وكلاهما يستخدم قاموساً \emph{قبلياً} (a priori) لتحفيز عمل منسق معقد. ومع ذلك، فإن الأنظمة التي يتحقق فيها التنسيق دون أي تبادل مباشر للمؤشرات منتشرة في الطبيعة والتكنولوجيا. في هذه المذكرة نقوم بإضفاء الطابع الرسمي على \textbf{YPSDC اللامركزي (d‑YPSDC)} – وهو بروتوكول تقوم فيه جميع الوكلاء بقراءة نفس المؤشر في وقت واحد من بيئة مشتركة، باستخدام قاموس موزع مسبقاً. لا يوجد أي اتصال بين الوكلاء، ومع ذلك يتم تحقيق كفاءة تنسيق \(K_{\mathrm{eff}} \gg 1\) من خلال التزامن في الوصول إلى المؤشر البيئي. نناقش تجليات d‑YPSDC في علم الأحياء (سلوك الأسراب)، والاقتصاد (رد فعل السوق على المؤشرات الاقتصادية الكلية)، وعلوم الحاسوب (مُزودات بيانات البلوكتشين)، والفيزياء الأساسية (اللاحتمية الكمية). يشير النموذج إلى وجود \textbf{حقل تنسيق} عميق يعمل كناقل عالمي للمؤشرات وينشط القاموس المشترك لقوانين الطبيعة.
	}
	
	\vspace{0.1cm}
	\noindent
	{\color{darkblue}\textbf{الكلمات المفتاحية:}} YPSDC، لامركزي، تنسيق، مؤشر بيئي، قاموس قبلي، حقل تنسيق، تشابك كمي، سلوك الأسراب، عقود ذكية.
	
	\vspace{0.5cm}
	\noindent
	\textcopyright\ 2026 أبحاث ياكوشيف. جميع الحقوق محفوظة.
	
	\tableofcontents
	\newpage
	
	\section{مقدمة: من YPSDC المركزي إلى بروتوكول لامركزي}
	
	يصف بروتوكول YPSDC الكلاسيكي تنسيق اثنين أو أكثر من الوكلاء، حيث يقوم أحدهم \textbf{بإرسال مؤشر قصير بنشاط}، بينما يقوم الآخرون \textbf{باستقباله} وتفعيل مدخل معقد من قاموس \emph{قبلي} مشترك \cite{YUCT}. في مثل هذا المخطط هناك دائماً مرسل مميز، وتكون طوبولوجيا التفاعل "نقطة‑إلى‑نقطة" أو "نجمية".
	
	تكشف ملاحظات الأنظمة البيولوجية والاجتماعية والفيزيائية عن فئة واسعة من الظواهر التي لا يصاحب التنسيق فيها تبادل مباشر للمؤشرات. الطيور في سرب تغير اتجاهها في وقت واحد دون قائد مرئي؛ السوق تتفاعل فوراً مع نشر رقم اقتصادي كلي؛ الجسيمات الكمية المتشابكة تظهر ارتباطات دون تبادل إشارات. في كل هذه الحالات يتحقق التنسيق لأن \textbf{جميع الوكلاء يقرؤون نفس المؤشر في نفس الوقت من البيئة المحيطة}، معتمدين على قاموس متفق عليه مسبقاً.
	
	الغرض من هذه المذكرة هو صياغة هذه الآلية رسمياً باسم \textbf{YPSDC اللامركزي (d‑YPSDC)}، ومناقشة بنيتها الرياضية، وإظهار أنها تعميم طبيعي للبروتوكول الكلاسيكي للأنظمة ذات طوبولوجيا "الناقل المشترك".
	
	\section{النموذج الرسمي لـ d‑YPSDC}
	
	\subsection{المكونات}
	
	يتكون النظام من المكونات التالية:
	
	\begin{itemize}[leftmargin=*]
		\item \(\mathcal{A} = \{A_1, A_2, \dots, A_N\}\) — مجموعة الوكلاء (مراقبون، خلايا، حيوانات، عقد شبكة).
		\item \(\mathcal{D}\) — \textbf{القاموس القبلي}، مشترك بين جميع الوكلاء وموزع خلال المرحلة غير المتصلة (أوفلاين). القاموس هو مجموعة من الأزواج \(\{ \kappa \to \text{فعل} \}\)، حيث \(\kappa\) مؤشر محتمل و\(\text{فعل}\) هو الإجراء المنسق.
		\item \(\mathcal{E}\) — فضاء حالات البيئة.
		\item \(S(t) \in \mathcal{E}\) — \textbf{المؤشر البيئي}، المتاح لجميع الوكلاء في وقت واحد. البيئة \emph{ليست} مرسلاً نشطاً؛ إنها توفر فقط حاملًا مادياً للمؤشر.
	\end{itemize}
	
	\subsection{البروتوكول}
	
	\begin{enumerate}[leftmargin=*, label=\textbf{المرحلة \arabic*:}]
		\item \textbf{توزيع القاموس غير المتصل (أوفلاين).} يتم إرسال القاموس الكامل \(\mathcal{D}\) إلى جميع الوكلاء بطريقة ما (وراثياً، ثقافياً، أو خلال خطوة تهيئة).
		\item \textbf{قراءة المؤشر المتصلة (أونلاين).} في الزمن \(t\) تولد البيئة المؤشر \(\kappa = S(t)\)، الذي يصل إلى جميع الوكلاء دون تأخير.
		\item \textbf{التفعيل.} يسترجع كل وكيل \(A_i\) بشكل مستقل المدخل \(\mathcal{D}(\kappa)\) من القاموس وينفذ الإجراء المحدد. لا يتم تبادل أي رسائل بين الوكلاء.
	\end{enumerate}
	
	\subsection{كفاءة التنسيق}
	
	على غرار YPSDC الكلاسيكي، تُعرَّف كفاءة التنسيق بأنها نسبة إنتروبيا الفعل المُطلَق إلى إنتروبيا المؤشر:
	\begin{equation}
		\Keff^{\mathrm{(d)}} = \frac{H(\text{الفعل الجماعي})}{H(\kappa)},
		\label{eq:Keff_d}
	\end{equation}
	حيث \(H(\text{الفعل الجماعي})\) هي إنتروبيا السلوك المتزامن للمجموعة بأكملها. ولأن المؤشر قصير جداً بينما يمكن أن يكون الفعل معقداً تعسفياً، فإن قيمة \(\Keff^{\mathrm{(d)}}\) يمكن أن تكون كبيرة جداً.
	
	\subsection{الفرق الطوبولوجي عن YPSDC الكلاسيكي}
	
	يظهر الفرق الجوهري بين البديلين في الشكل~\ref{fig:topology}.
	
	\begin{figure}[htbp]
		\begin{otherlanguage}{english}
			\centering
			\begin{tikzpicture}[
				agent/.style={circle, draw=blue!70, fill=blue!15, minimum size=1.2cm, font=\small},
				dict/.style={rectangle, draw=green!50!black, fill=green!10, rounded corners, minimum width=3cm, minimum height=0.9cm, font=\small},
				env/.style={rectangle, draw=orange!70, fill=orange!10, rounded corners, minimum width=3cm, minimum height=0.9cm, font=\small},
				arrow/.style={->, >=stealth, thick},
				dashedarrow/.style={->, >=stealth, thick, dashed},
				]
				
				% ---- YPSDC الكلاسيكي (أعلى) ----
				\begin{scope}[local bounding box=classic]
					\node[dict] (dict1) at (0,0) {القاموس \(\mathcal{D}\)};
					\node[agent] (A1) at (-3, -2) {الوكيل A};
					\node[agent] (A2) at (3, -2) {الوكيل B};
					\draw[dashedarrow] (dict1) -- (A1) node[midway, left] {غير متصل};
					\draw[dashedarrow] (dict1) -- (A2) node[midway, right] {غير متصل};
					\draw[arrow] (A1) -- (A2) node[midway, above] {المؤشر \(\kappa\)};
				\end{scope}
				\node[above=0.2cm of classic.north, font=\bfseries] {YPSDC الكلاسيكي (طوبولوجيا مركزية)};
				
				% ---- d‑YPSDC (أسفل) ----
				\begin{scope}[local bounding box=decentr, yshift=-6cm]
					\node[env] (env) at (0,1) {البيئة \(S(t)\)};
					\node[dict] (dict2) at (0,-0.5) {القاموس \(\mathcal{D}\)};
					\node[agent] (B1) at (-4, -2.5) {الوكيل 1};
					\node[agent] (B2) at (0, -2.5) {الوكيل 2};
					\node[agent] (B3) at (4, -2.5) {الوكيل 3};
					\draw[arrow] (env) -- (B1) node[midway, left] {\(\kappa\)};
					\draw[arrow] (env) -- (B2) node[midway, right] {\(\kappa\)};
					\draw[arrow] (env) -- (B3) node[midway, right] {\(\kappa\)};
					\draw[dashedarrow] (dict2) -- (B1);
					\draw[dashedarrow] (dict2) -- (B2);
					\draw[dashedarrow] (dict2) -- (B3);
					% تأشير واضح لغياب الوصلات بين الوكلاء
					\draw[red, thick, dotted] (B1) -- (B2) node[midway, above, red] {لا رابط};
					\draw[red, thick, dotted] (B2) -- (B3) node[midway, above, red] {لا رابط};
				\end{scope}
				\node[above=0.2cm of decentr.north, font=\bfseries] {d‑YPSDC (طوبولوجيا لامركزية)};
				
			\end{tikzpicture}
		\end{otherlanguage}
		\caption{مقارنة الطوبولوجيات. في YPSDC الكلاسيكي، يرسل الوكيل A مؤشراً بنشاط إلى الوكيل B. في d‑YPSDC اللامركزي، يقرأ جميع الوكلاء نفس المؤشر من البيئة ولا يتبادلون أي رسائل.}
		\label{fig:topology}
	\end{figure}
	
	\section{أمثلة على d‑YPSDC في أنظمة مختلفة}
	
	\subsection{علم الأحياء: سلوك الأسراب و"الموافقة الضمنية"}
	
	لوحظت أفعال جماعية متزامنة (تحولات مفاجئة لسرب من الأسماك، إقلاع متزامن لسرب طيور، تجمد السنجاب الأرضي) في العديد من أنواع الحيوانات دون إشارات صوتية أو بصرية خاصة بكل فرد. بدلاً من ذلك، تستشعر جميع الأفراد \textbf{في آن واحد} تغيراً في معامل بيئي – هبوب ريح، ظل مفترس، تغير في الإضاءة – وتتفاعل وفقاً لقاموس فطري (جيني).
	
	مثال الشعر البشري الذي نوقش سابقاً يتناسب أيضاً مع هذا المخطط: الشعر الطويل على الرأس يمكن أن يعمل ككاشف حساس جداً لتيارات الهواء ودرجة الحرارة، مما يمكن مجموعة من الناس من تسجيل تغير الطقس بشكل متزامن دون تبادل الكلمات.
	
	\subsection{الاقتصاد والمالية}
	
	نشر مؤشر اقتصادي رئيسي (على سبيل المثال، معدل البطالة في الولايات المتحدة) يثير رد فعل فوري ومتماسك من ملايين المشاركين في السوق. التجار لا يتواصلون مع بعضهم البعض؛ يستخدمون "قاموساً" مشتركاً – خوارزميات تداول، نماذج اقتصادية، أو ببساطة توقعات تكونت من الخبرة السابقة. المؤشر (القيمة العددية للمؤشر) تقدمه البيئة (محطات المعلومات) لجميعهم في وقت واحد، ويتحرك السوق ككيان واحد.
	
	\subsection{البلوكتشين والعقود الذكية}
	
	تعتمد التطبيقات اللامركزية (dApps) غالباً على \textbf{مُزودات البيانات (Oracles)} – مصادر بيانات خارجية تزود البلوكتشين بمعلومات عن العالم الحقيقي. يلعب العقد الذكي دور القاموس الموزع على جميع عقد الشبكة. عندما ينشر المُزود مؤشراً (على سبيل المثال، سعر ETH/USD في الزمن \(t\))، تقوم جميع العقد بتفعيل الفرع المقابل من العقد في وقت واحد دون توافق إضافي. هذا هو تطبيق خالص لـ d‑YPSDC في بيئة رقمية.
	
	\subsection{التشابك الكمي}
	
	ربما أعمق تحقق فيزيائي لـ d‑YPSDC يحدث في ميكانيكا الكم. تأمل زوجاً من الجسيمات المتشابكة. دالتها الموجية المشتركة هي قاموس يحدد الارتباطات بين نتائج القياس المحتملة. عندما يتم قياس أحد الجسيمات، يمكن اعتبار نتيجته \textbf{مؤشراً بيئياً} يمكن للجسيم الثاني الوصول إليه في نفس الوقت بفضل الخصائص اللاحتمالية للحقل الكمي. الجسيم الثاني "يفعل" الحالة المرتبطة دون تلقي أي إشارة كلاسيكية. في إطار d‑YPSDC، يُفسر هذا على أنه قراءة متزامنة لنفس المؤشر من وسط فيزيائي مشترك (الحقل الكمي)، دون حاجة لنقل معلومات بين الجسيمات.
	
	\section{الأساس الوجودي: المادة كتنسيق، والهندسة كنتيجة}
	
	الفرق الرئيسي بين YUCT ونظريات الحقل القياسية هو أن \textbf{حقل التنسيق لا يوجد في زمكان فارغ}. على العكس، الزمكان نفسه هو تجلٍ ثانوي وناشئ لواقع أكثر جوهرية – شبكة من العناصر المنسقة. هذا الافتراض له نتيجتان رئيسيتان تنطبقان مباشرة على YPSDC اللامركزي.
	
	\textbf{1. المادة = تنسيق.} مفهوم "الجسيم الأولي" أو "الوكيل" ليس مستقلاً عن عملية التنسيق. العنصر يوجد فقط بقدر ما يشارك في أفعال التنسيق ويمكن تعريفه من خلال قاموسه. خارج شبكة التنسيق لا توجد مادة؛ تبقى فقط نقاط رياضية شكلية بلا محتوى فيزيائي. في سياق d‑YPSDC هذا يعني أن الوكلاء الذين يقرؤون المؤشر البيئي ليسوا مراقبين سلبيين – وجودهم المجرد كوحدات تتصرف بشكل متماسك \emph{يخلق} حقل التنسيق \(\Psi_{MN}\) الذي يوفر بعد ذلك المؤشرات لهم.
	
	\textbf{2. الهندسة لا وجود لها دون مادة.} مفارقة القرص الدوار (المعروفة في الأدبيات بمفارقة إهرنفست) تبين أن الهندسة المكانية لا يمكن أن تُعطى \emph{قبلياً} بل تحددها حالة الحركة والقوى الترابطية الداخلية للجسم المادي. في YUCT يُعمم هذا إلى القول بأن المقياس الزمكاني \(g_{\mu\nu}(x)\) هو كمية \emph{فعالة} تنشأ من متوسط تفاعلات التنسيق. حيث يغيب التنسيق (\(K_{\mathrm{eff}} \to 1\))، تتدهور الهندسة وتفقد أهميتها التشغيلية. في d‑YPSDC يتجلى هذا في حقيقة أن "الموصلية" الفعالة للبيئة للمؤشر تعتمد على كثافة وترابط الوكلاء المنسقين، وليس على خصائص الفضاء المطلق.
	
	وهكذا، \textbf{البيئة التي تسلم المؤشر ليست وعاءً سلبياً بل هي حقل التنسيق نفسه، الناشئ عن الوكلاء الماديين}. إن وجود وكلاء بقاموس مشترك هو الذي "يشغل" الحقل \(\Psi_{MN}\) ويمنحه البنية التي تمكن من تسليم المؤشرات المتزامنة. في غياب الوكلاء (أو عندما يُدمر القاموس) يعود الحقل إلى حالة متماثلة غير منسقة، تعادل غياب الزمكان الكلاسيكي. وهذا يتوافق تماماً مع الصورة التي يكون فيها المكان والزمان والمادة ثلاثة جوانب لا تنفصل لعملية تنسيق واحدة.
	
	\section{النموذج الرياضي لـ d‑YPSDC في الفضاء 19‑البعدي لـ YUCT}
	\label{sec:math-model}
	
	\subsection{حقل التنسيق 19‑البعدي وفعله}
	
	الساحة الأساسية لـ YUCT هي متعددة شعب شبه ريمانية 19‑بعدية \(\mathcal{M}_{19}\) بمترية \(G_{MN}\) ذات إشارة \((+,-,\dots,-)\). هي مفسرة على الزمكان رباعي الأبعاد \(M_4\):
	\begin{equation}
		\mathcal{M}_{19} = M_4 \times F_{15},
	\end{equation}
	حيث \(F_{15}\) فضاء داخلي مدمج لدرجات حرية التنسيق. الحقل المحوري هو \textbf{حقل التنسيق} \(\Psi_{MN}(X)\)، الذي يشفر القاموس الموزع والوسط الناقل للمؤشر. أبسط فعل له الشكل
	\begin{equation}
		S_{19} = \frac{1}{2\kappa_{19}} \int d^{19}X \sqrt{-G}\, R(G) + \int d^{19}X \sqrt{-G}\, \mathcal{L}_{\mathrm{coord}},
		\label{eq:S19}
	\end{equation}
	حيث يوفر الانحناء \(R(G)\) الهندسة الخلفية، ولامرانج التنسيق هو
	\begin{equation}
		\mathcal{L}_{\mathrm{coord}} = -\frac{1}{4} \Tr(\Psi_{MN}\Psi^{MN}) + \frac{1}{2} G^{MN}\,\partial_M \phi\,\partial_N \phi - V(\phi)
		\label{eq:Lcoord}
	\end{equation}
	ويحتوي على حقل القياس \(\Psi_{MN}\) (شبيه بقوة تفاعل التنسيق) وحقل قياسي \(\phi = \ln(K_{\mathrm{eff}}/K_0)\) يحدد كفاءة التنسيق المحلية.
	
	\subsection{إدخال الوكلاء والقاموس}
	
	يُمثل الوكلاء \(A_i\)، \(i=1,\dots,N\)، كمصادر نقطية متمركزة في \(M_4\) و"منبسطة" على \(F_{15}\) وفقاً لشكل قاموسهم. كل وكيل يمتلك \textbf{دالة قاموس} \(f_{\mathcal{D}}^{(i)}(\kappa)\) تربط المؤشر \(\kappa \in \mathcal{E}\) بالفعل المطلوب. يُعطى تفاعل الوكلاء مع الحقل \(\Psi_{MN}\) بالحد
	\begin{equation}
		S_{\mathrm{int}} = \sum_{i=1}^{N} \int d\tau_i \left[ \frac{1}{2} m_i \dot a_i^2 - \lambda_i \bigl( a_i - f_{\mathcal{D}}^{(i)}(\kappa) \bigr)^2 \right],
		\label{eq:Sint}
	\end{equation}
	حيث \(a_i\) هو الفعل المنفذ فعلياً و\(\tau_i\) هو الزمن الذاتي للوكيل. المؤشر \(\kappa\) هنا ليس معاملاً خارجياً، بل قيمة الحقل \(\Psi_{MN}\) (أو إسقاطه على \(M_4\)) في موقع الوكيل:
	\begin{equation}
		\kappa(x_i) = \int_{F_{15}} d^{15}Y \sqrt{-G^{(15)}}\, \mathcal{O}[\Psi_{MN}](x_i, Y),
		\label{eq:kappa}
	\end{equation}
	حيث \(\mathcal{O}\) مؤثر يسقط حقل التنسيق على فضاء المؤشر. في أبسط الحالات \(\mathcal{O}[\Psi] = \Tr(\Psi)\) أو ببساطة \(\phi\).
	
	\subsection{معادلات الحركة وتسليم المؤشر}
	
	تغير الفعل الكلي \(S_{19} + S_{\mathrm{int}}\) بالنسبة لـ \(\Psi_{MN}\) يعطي المعادلة
	\begin{equation}
		D_M \Psi^{MN} = J^{N}_{\mathrm{coord}},
		\label{eq:Psi-eq}
	\end{equation}
	حيث \(J^{N}_{\mathrm{coord}}\) هو تيار التنسيق الكلي للوكلاء:
	\begin{equation}
		J^{N}_{\mathrm{coord}}(X) = \sum_{i=1}^{N} \int d\tau_i\, \frac{\delta^{(19)}(X - X_i(\tau_i))}{\sqrt{-G}}\, \frac{\partial \mathcal{L}_{\mathrm{int}}}{\partial (\partial_N \Psi)}.
	\end{equation}
	للحقول المتغيرة ببطء ولقاموس ثابت، تختزل المعادلة الفعالة إلى معادلة موجية على \(M_4\) بمصادر:
	\begin{equation}
		\Box \phi - m_{\mathrm{eff}}^2 \phi = \sum_{i} g_i \delta^{(4)}(x - x_i),
		\label{eq:phi-eq}
	\end{equation}
	حيث \(m_{\mathrm{eff}}^2 = V''(0)\)، وثوابت الارتباط \(g_i\) متناسبة مع "جهارة" المؤشر.
	
	حل المعادلة~(\ref{eq:phi-eq}) في التقريب شبه الساكن هو
	\begin{equation}
		\phi(x) = \sum_i g_i \, G_{\mathrm{ret}}(x, x_i),
	\end{equation}
	مع \(G_{\mathrm{ret}}\) دالة جرين المتأخرة لحقل قياسي ذي كتلة. إذا كان جميع الوكلاء موجودين في منطقة أبعادها أصغر بكثير من طول كومبتون \(\lambda_c = 1/m_{\mathrm{eff}}\)، فإن \(\phi(x)\) عملياً متطابق عند جميع النقاط \(x_i\). هذا يفسر \textbf{القراءة المتزامنة لنفس المؤشر} بواسطة جميع الوكلاء: المؤشر، الناشئ عن البيئة (أي بشروط حدودية أو خلفية كونية)، ينتشر عبر \(M_4\) ويظهر، بسبب طوله الموجي الطويل، متطابقاً لجميع المراقبين المحليين.
	
	وبالتالي، فإن تزامن التفعيل لا يتطلب إشارات أسرع من الضوء: إنه نتيجة طبيعية لحقيقة أن حقل التنسيق \(\phi\) يعمل كخلفية متجانسة على مقياس المجموعة. في نفس الوقت، كل وكيل يتفاعل حصراً مع الحقل، وليس مع وكلاء آخرين؛ لا يحدث أي تبادل مباشر للمعلومات – وهذا هو جوهر YPSDC اللامركزي.
	
	\subsection{الارتباط بكفاءة التنسيق \(K_{\mathrm{eff}}\)}
	
	تُعرَّف كفاءة التنسيق المحلية بأنها نسبة السعة المعلوماتية للقاموس إلى طول المؤشر. بدلالة الحقل \(\phi\):
	\begin{equation}
		\Keff(x) = K_0 e^{\phi(x)},
		\label{eq:Keff-field}
	\end{equation}
	حيث \(K_0 \approx 1\) هي قيمة الخلفية في الفراغ (غياب وكلاء منسقين). نقل مؤشر بيئي يخلق اضطراباً \(\delta\phi\)، بحيث تكون في المنطقة التي يشغلها الوكلاء
	\begin{equation}
		\Keff^{\mathrm{(d)}} = K_0 \exp\!\left( \sum_i g_i G_{\mathrm{ret}}(x, x_i) \right).
	\end{equation}
	لمجموعة من \(N\) وكيل متطابق يقعون ضمن حجم \(V \ll \lambda_c^3\)،
	\begin{equation}
		\Keff^{\mathrm{(d)}} \approx K_0 \exp\!\left( \frac{N g}{4\pi} \frac{e^{-m_{\mathrm{eff}} r_0}}{r_0} \right),
	\end{equation}
	حيث \(r_0\) هي المسافة المميزة بين الوكلاء. بما أن \(N g > 0\) (المؤشر يحسن التنسيق)، فإن قيمة \(\Keff^{\mathrm{(d)}}\) يمكن أن تتجاوز الواحد بكثير للمجموعات الكبيرة \(N\) أو الارتباط القوي \(g\).
	
	هذا يظهر مباشرة أن \textbf{YPSDC اللامركزي يقدم \(K_{\mathrm{eff}} \gg 1\) دون أي تبادل للمؤشرات داخل المجموعة}. نمو الكفاءة يعود حصراً لهندسة الحقل \(\phi\): جميع الوكلاء مغمورون في حقل تنسيق مشترك، والذي، عند إثارته بمؤشر خارجي، ينشط قواميسهم في نفس الوقت.
	
	\section{نظامان لنفس البروتوكول: YPSDC مركزي مقابل لامركزي}
	\label{sec:two-regimes}
	
	الأمثلة السابقة والنموذج الرياضي قد تعطي انطباعاً بأن d‑YPSDC هو بروتوكول تنسيق جديد جوهرياً. ليس كذلك. وجودياً هناك \textbf{آلية YPSDC واحدة فقط}: حقل التنسيق \(\Psi_{MN}\) ينقل المؤشرات، وجميع الوكلاء يتفاعلون حصراً مع هذا الحقل، مفعلين مدخلات قاموس \emph{قبلي} مشترك. التمييز بين YPSDC "المركزي" و d‑YPSDC "اللامركزي" هو ظاهري بحت ويعكس المصدر الغالب لتيار التنسيق \(J^{N}_{\mathrm{coord}}\) الذي يظهر في معادلات الحقل (انظر الملحق~أ من \cite{YUCT}).
	
	\subsection{النظامان}
	\begin{enumerate}[leftmargin=*, label=\textbf{(\arabic*)}]
		\item \textbf{النظام المركزي (YPSDC الكلاسيكي).} التيار \(J^{N}_{\mathrm{coord}}\) يهيمن عليه وكيل متمركز (المرسل). هذا الوكيل يثير الحقل بنشاط، خالقاً اضطراباً منتشراً يصل إلى واحد أو أكثر من المستقبلين المحددين. هذا النظام يصف التواصل نقطة‑إلى‑نقطة أو طوبولوجيا النجمة (مثل برج GSM يبث إلى هواتف محمولة، أو إشارة توقيت GMT).
		
		\item \textbf{النظام اللامركزي (d‑YPSDC).} التيارات المتمركزة للوكلاء الأفراد مهملة مقارنة بخلفية كونية متغيرة ببطء – نمط بيئي لـ \(\Psi_{MN}\). الخلفية تقدم مؤشراً متاحاً لجميع الوكلاء في وقت واحد. هذا النظام يصف التشابك الكمي، سلوك الأسراب، استشعار النصاب (quorum sensing)، رد فعل الأسواق المالية على المؤشرات الاقتصادية الكلية، وتفعيل العقود الذكية بواسطة مُزودات البلوكتشين.
	\end{enumerate}
	
	\subsection{الانتقال بين النظامين}
	المعامل اللامبعدي
	\[
	\Theta = \frac{|J_{\mathrm{وكيل}}|}{|J_{\mathrm{بيئة}}|},
	\]
	حيث \(J_{\mathrm{وكيل}}\) هو التيار الناشئ عن أقوى مرسل فردي و\(J_{\mathrm{بيئة}}\) هو التيار الفعال للخلفية البيئية، يتحكم في الانتقال بين النظامين:
	\[
	\begin{cases}
		\Theta \gg 1 & \text{نظام مركزي (YPSDC كلاسيكي)}, \\
		\Theta \ll 1 & \text{نظام لامركزي (d‑YPSDC)}, \\
		\Theta \sim 1 & \text{سلوك مختلط}.
	\end{cases}
	\]
	
	\subsection{الوضع الوجودي}
	d‑YPSDC لا يحتاج \textbf{إلى حقول إضافية، ولا تعديل للاغرانجي، ولا مراجعة للمبادئ الوجودية} لـ YUCT. الفعل 19‑البعدي \(S_{19}\) وحقول المراقب \(\phi_1,\phi_2\) (الملحق~أ) تحتوي بالفعل على كل ما هو ضروري لكلا النظامين. مصطلح "d‑YPSDC" يُحتفظ به فقط كتسمية ظاهرية مريحة للفئة الواسعة من الظواهر التي يتحقق فيها التنسيق دون مرسل نشط يمكن تحديده.
	
	هذا التوضيح ضروري للاكتمال المنطقي لإطار YUCT: النظرية تصف جميع أشكال التنسيق – من التواصل الإنساني الصريح إلى الارتباطات اللاحتمية لميكانيكا الكم – من خلال \textbf{بروتوكول واحد} يُطبق على \textbf{حقل تنسيق واحد}.
	
	\subsection{مثال: التشابك الكمي كـ d‑YPSDC}
	
	في الحالة الكمية، القاموس هو الدالة الموجية المشتركة \(\Psi_{AB}\) لزوج متشابك. اختر المؤثر \(\mathcal{O}\) في (\ref{eq:kappa}) بحيث يتطابق المؤشر \(\kappa\) مع نتيجة قياس على الجسيم الأول. عندها تصف المعادلة~(\ref{eq:phi-eq}) انتشار هذه النتيجة (عبر حقل التنسيق \(\phi\)) إلى الجسيم الثاني. وبما أن الحقل ذو كتلة وسكوني، فإن الارتباط ينشأ فوراً بمعنى التفعيل المتزامن، ولكن بدون نقل طاقة. هذا هو التفسير الهندسي لانتهاك متباينة بل في YUCT.
	
	رياضياً، درجة الانتهاك \(S\) تُعبر عنها بالكفاءة:
	\begin{equation}
		S = 2\sqrt{2} \frac{\Keff^{\mathrm{(d)}}}{\Keff^{\mathrm{(d)}} + 1},
	\end{equation}
	بحيث تصل القيمة الكمية القصوى \(2\sqrt{2}\) عندما \(\Keff^{\mathrm{(d)}} \to \infty\) (قاموس مثالي)، بينما تختفي الارتباطات عندما \(\Keff^{\mathrm{(d)}} \to 1\) (تدمير القاموس).
	
	هذا النموذج يوحد علم الأحياء والاقتصاد وفيزياء الكم ضمن المعادلات المشتركة (\ref{eq:Psi-eq})–(\ref{eq:Keff-field})، مبرهناً على عالمية حقل التنسيق \(\Psi_{MN}\) وبروتوكول d‑YPSDC.
	
	\section{حقل التنسيق كناقل عالمي للمؤشرات}
	
	إذا كان d‑YPSDC صالحاً للأنظمة الكمية والمجموعات البيولوجية والشبكات الاجتماعية، فمن الطبيعي افتراض وجود \textbf{حقل تنسيق} موحد – وسط عالمي:
	
	\begin{itemize}
		\item يخزن "القاموس العالمي" لقوانين الطبيعة (تماثلات، ثوابت، معادلات حركة)؛
		\item يوصل المؤشرات (نتائج القياس، التقلبات، المؤثرات الخارجية) إلى جميع عناصر النظام في وقت واحد؛
		\item يضمن تزامن التفعيل الذي يُلاحظ على أنه لاحتمية كمية، سلوك جماعي، أو استجابة سريعة للسوق.
	\end{itemize}
	
	في هذا السياق، يتبين أن التواصل الكلاسيكي (نقل إشارة من وكيل إلى آخر) هو \textbf{حالة خاصة} تظهر عندما لا تكون البيئة قادرة على إيصال المؤشر إلى جميع الوكلاء في وقت واحد (على سبيل المثال، بسبب مستوى ضوضاء مرتفع)، مما يجبر الوكلاء على مضاعفة المؤشر عن طريق التواصل المباشر.
	
	هذه الفرضية لا تتعارض مع النسبية، لأنه لا تنتقل طاقة أو معلومات أسرع من الضوء: ينتشر المؤشر عبر البيئة بسرعة لا تتجاوز \(c\) ويصبح متاحاً لجميع الوكلاء فقط بعد وصوله إلى جوارهم المحلي. تزامن التفعيل يعود حصراً للقاموس المشترك وهوية المؤشر، وليس للإشارات الأسرع من الضوء.
	
	\section{الاستنتاج}
	
	قدمنا مفهوم YPSDC اللامركزي (d‑YPSDC) – بروتوكول يتم فيه التنسيق من خلال القراءة المتزامنة لمؤشر بيئي مشترك من قبل جميع الوكلاء باستخدام قاموس \emph{قبلي} موزع مسبقاً. هذا المخطط يفسر مجموعة واسعة من الظواهر – من التزامن البيولوجي إلى اللاحتمية الكمية – دون اللجوء إلى أي تواصل خفي.
	
	يقود d‑YPSDC بشكل طبيعي إلى مفهوم \textbf{حقل تنسيق} ككيان أساسي مسؤول عن الاتساق الكوني للعمليات الطبيعية. النموذج الرياضي القائم على الفعل 19‑البعدي لـ YUCT يظهر أن التنسيق اللامركزي ينشأ من معادلات موحدة للحقل \(\Psi_{MN}\)، وأن الكفاءة \(K_{\mathrm{eff}}\) تنمو أسياً مع عدد الوكلاء وقوة الارتباط. مزيد من الصياغة الرسمية لهذا الحقل وربطه بنظرية الكم والنسبية العامة تشكل برنامجاً للبحث المستقبلي.
	
	\begin{thebibliography}{9}
		\begin{otherlanguage}{english}
			\bibitem{YUCT} A.~V.~Yakushev, \emph{Yakushev Unified Coordination Theory (YUCT)}, Zenodo, DOI:10.5281/zenodo.18444599 (2026).
			\bibitem{AppendixB} A.~V.~Yakushev, Appendix B: Temporal Coordination Paradox, in \emph{YUCT} (2026).
			\bibitem{Blockchain} S.~Nakamoto, ``Bitcoin: A Peer-to-Peer Electronic Cash System'' (2008).
			\bibitem{Ben-Jacob} E.~Ben-Jacob, ``Bacterial self-organization: co-enhancement of complexification and adaptability'', \emph{Physica A} (2003).
			\bibitem{EPR} A.~Einstein, B.~Podolsky, N.~Rosen, ``Can Quantum-Mechanical Description of Physical Reality Be Considered Complete?'', \emph{Phys. Rev.} \textbf{47}, 777 (1935).
		\end{otherlanguage}
	\end{thebibliography}
	
\end{document}